\question[5]
Eine Messstation speichert ihre Messwerte in einer Zeile, getrennt durch Semikolons:
\texttt{zeile = '12;7;30;5'}

Schreibe Python-Anweisungen, die
\begin{itemize}
\item die Zeile in die einzelnen Werte zerlegen und daraus eine Liste \texttt{werte} mit ganzen Zahlen erzeugen,
\item die Summe und den größten Wert ausgeben,
\item die Werte mit \texttt{' + '} verbunden ausgeben, also \texttt{12 + 7 + 30 + 5}.
\end{itemize}

\begin{solutionbox}{7cm}
\begin{lstlisting}
zeile = '12;7;30;5'
teile = zeile.split(';')      # ['12', '7', '30', '5']
werte = []
for t in teile:
    werte.append(int(t))
print(sum(werte), max(werte))     # 54 30

texte = []
for w in werte:
    texte.append(str(w))
print(' + '.join(texte))
\end{lstlisting}
(Für die letzte Ausgabe geht auch \texttt{print(' + '.join(teile))}, denn \texttt{teile} enthält schon Strings.)
\end{solutionbox}
